% 编译提示：使用 XeLaTeX 编译
% 需要 ctex 宏包支持中文

\documentclass[a4paper,9pt]{article}
\usepackage{amsmath}
\usepackage{amssymb}
\usepackage{geometry}
\usepackage{ctex}
\usepackage{xcolor}
\usepackage{multicol}
\usepackage{titlesec}
\usepackage{fancyhdr}
\usepackage{tcolorbox}
\usepackage{graphicx}
\usepackage[version=4]{mhchem}

\geometry{left=1cm,right=1cm,top=1cm,bottom=1.8cm}
\setlength{\columnsep}{0.8cm}
\setlength{\columnseprule}{0.4pt}
\renewcommand{\columnseprulecolor}{\color{gray!50}}

\definecolor{sectioncolor}{RGB}{0,102,204}
\definecolor{subsectioncolor}{RGB}{0,153,153}
\definecolor{keycolor}{RGB}{204,102,0}
\definecolor{derivcolor}{RGB}{100,149,237}
\definecolor{boxbg}{RGB}{240,248,255}
\definecolor{keybg}{RGB}{255,248,240}

\titleformat{\section}
  {\normalfont\large\bfseries\color{sectioncolor}}
  {\thesection}{0.5em}{}
\titlespacing*{\section}{0pt}{1ex plus 0.5ex minus 0.2ex}{0.8ex plus 0.2ex}

\titleformat{\subsection}
  {\normalfont\normalsize\bfseries\color{subsectioncolor}}
  {\thesubsection}{0.5em}{}
\titlespacing*{\subsection}{0pt}{0.8ex plus 0.3ex minus 0.1ex}{0.5ex plus 0.1ex}

\setcounter{secnumdepth}{4}
\setcounter{tocdepth}{4}

\newenvironment{keypoint}
{\par\vspace{0.3em}\noindent\begin{tcolorbox}[
  colback=keybg,
  colframe=keycolor,
  boxrule=0.5pt,
  arc=2pt,
  left=2pt,
  right=2pt,
  top=2pt,
  bottom=2pt,
  boxsep=0pt,
  ]
\noindent\textcolor{keycolor}{\textbf{关键点：}}}
{\end{tcolorbox}\par\vspace{0.3em}}

\newenvironment{derivation}
{\par\vspace{0.3em}\noindent\begin{tcolorbox}[
  colback=boxbg,
  colframe=derivcolor,
  boxrule=0.5pt,
  arc=2pt,
  left=2pt,
  right=2pt,
  top=2pt,
  bottom=2pt,
  boxsep=0pt,
  ]
\scriptsize\noindent\textcolor{derivcolor}{\textbf{推导：}}\par}
{\end{tcolorbox}\par\vspace{0.3em}}

\newenvironment{examplebox}
{\par\vspace{0.3em}\noindent\begin{tcolorbox}[
  colback=boxbg,
  colframe=derivcolor,
  boxrule=0.5pt,
  arc=2pt,
  left=2pt,
  right=2pt,
  top=2pt,
  bottom=2pt,
  boxsep=0pt,
  ]
\scriptsize\noindent\textcolor{derivcolor}{\textbf{示例：}}\par
\setlength{\parskip}{0.1ex}\setlength{\itemsep}{0pt}\setlength{\parsep}{0pt}}
{\end{tcolorbox}\par\vspace{0.3em}}

\pagestyle{fancy}
\fancyhf{}
\fancyfoot[C]{\scriptsize3-\thepage}
\renewcommand{\headrulewidth}{0pt}
\renewcommand{\footrulewidth}{0.4pt}
\renewcommand{\footrule}{\hbox to\headwidth{\color{gray!50}\leaders\hrule height \footrulewidth\hfill}}

\title{\vspace{-2em}\Large\bfseries\color{sectioncolor} Chapter 3 Collections - 配分函数的具体求值（Lect.4）\vspace{-1em}}
\author{}
\date{\today}

\begin{document}
\maketitle
\thispagestyle{fancy}

\begin{multicols}{2}

\scriptsize{\tableofcontents}

	\section{如何真正去“算” $q$}

	Lect.~3 解决了“知道 $q$ 以后怎么算 $A,U,S$”；Lect.~4 解决的是“给定某组能级，如何实际求出 $q$”。

	对一组能级写出
	\[
		q=\sum_i e^{-\varepsilon_i/kT}
	\]
	之后，第一判断永远是：\textbf{很多项参与，还是只有少数项参与？}

	\begin{keypoint}
		总原则：
		\begin{itemize}
			\item 若能级间隔 $\ll kT$，很多项有贡献，常可把和式近似为积分；
			\item 若能级间隔 $\gg kT$，只需保留前几项，积分近似很差；
			\item 若能级间隔 $\sim kT$，逐项算到项变小为止。
		\end{itemize}
	\end{keypoint}

	\subsection{$kT$ 判据的等间距图像}
	若能级等间距为
	\[
		0,\Delta,2\Delta,3\Delta,\dots
	\]
	则
	\[
		q'=1+e^{-\Delta/kT}+e^{-2\Delta/kT}+e^{-3\Delta/kT}+\cdots
	\]
	可分三类：
	\begin{itemize}
		\item $\Delta\gg kT$：第二项都很小，$q'\approx1$；
		\item $\Delta\ll kT$：很多项贡献，$q'\gg1$；
		\item $\Delta\sim kT$：只有少数几项有贡献。
	\end{itemize}

	\begin{keypoint}
		只有满足
		\[
			\varepsilon_i\lesssim kT
		\]
		的能级才显著贡献到配分函数，因此 $q$ 也可以理解成“可及态数”的量度。
	\end{keypoint}

	\subsection{积分近似什么时候可用}
	若很多项贡献，和式可近似成积分。例如上式可写成
	\[
		q'\approx \int_0^{\infty} e^{-i\Delta/kT}\,di.
	\]
	若只有 1--2 项重要，则积分会严重高估，必须逐项求和。

	此外，当只能逐项算时，导数也可逐项算：
	\[
		\left(\frac{\partial q}{\partial T}\right)_V
		=\frac1{kT^2}\sum_i \varepsilon_i e^{-\varepsilon_i/kT}.
	\]

	\section{degeneracy：对态求和还是对能级求和}

	若不同波函数具有相同能量，则系统有简并。此时配分函数本质上仍是对\textbf{态}求和，只是更方便改写成对\textbf{能级}求和：
	\[
		q=\sum_j g_j e^{-\varepsilon_j/kT},
	\]
	其中 $g_j$ 为第 $j$ 个能级的简并度。

	\begin{examplebox}
		二维正方形箱中
		\[
			E_{n_x,n_y}=\frac{h^2}{8ma^2}(n_x^2+n_y^2)
		\]
		若 $n_x\neq n_y$，则 $(n_x,n_y)$ 与 $(n_y,n_x)$ 为两个不同波函数但能量相同，因此存在二重简并。
	\end{examplebox}

	\begin{keypoint}
		见到题目给出“某能级的 degeneracy 为 $g_j$”，第一反应就是
		\[
			q=\sum_j g_j e^{-\varepsilon_j/kT}.
		\]
		这和 lect1 的布居公式
		\[
			n_j=\frac{N}{q}g_j e^{-\varepsilon_j/kT}
		\]
		是一一对应的。
	\end{keypoint}

	\section{平动配分函数}

	\subsection{粒子在箱中模型}
	一维箱中粒子能级：
	\[
		E_n=\frac{n^2h^2}{8ma^2},\qquad n=1,2,3,\dots
	\]
	其中 $a$ 是箱长，$m$ 是粒子质量。

	先估计有多少项贡献：令 $E_{n_{\max}}\sim kT$，则
	\[
		\frac{n_{\max}^2h^2}{8ma^2}=kT
		\quad\Rightarrow\quad
		n_{\max}^2=\frac{8ma^2kT}{h^2}.
	\]
	常温下气体平动能级极密，因此 $n_{\max}$ 极大，积分近似完全安全。

	\subsection{一维结果}
	于是
	\[
		q_{\mathrm{trans},1D}=\sum_{n=1}^{\infty} e^{-E_n/kT}
		\approx \int_0^{\infty} e^{-E_n/kT}\,dn.
	\]
	这里把下限从 $1$ 改成 $0$，是因为常温气体满足 $n_{\max}\gg 1$；多加的那个“假想 $n=0$ 项”相对成千上万真实可及态可以忽略，不会改变积分结果。
	借助高斯积分得
	\[
		q_{\mathrm{trans},1D}=\left(\frac{2\pi mkT}{h^2}\right)^{1/2}a.
	\]

	定义热波长
	\[
		\Lambda=\left(\frac{h^2}{2\pi mkT}\right)^{1/2},
	\]
	则
	\[
		q_{\mathrm{trans},1D}=\frac{a}{\Lambda}.
	\]

	\subsection{三维结果}
	三维箱中能级为
	\[
		E_{n_x,n_y,n_z}=\frac{h^2}{8ma^2}(n_x^2+n_y^2+n_z^2).
	\]
	由于三方向独立，配分函数可乘，得到
	\[
		q_{\mathrm{trans},3D}=\left(\frac{2\pi mkT}{h^2}\right)^{3/2}V
		=\frac{V}{\Lambda^3}.
	\]

	\begin{keypoint}
		平动公式是本章最高频：
		\[
			q_{\mathrm{trans},3D}=\left(\frac{2\pi mkT}{h^2}\right)^{3/2}V=\frac{V}{\Lambda^3}.
		\]
		平动几乎总在“高温极限”，所以最常见做法就是直接上积分结果。
	\end{keypoint}

	\begin{examplebox}
		对应习题：
		\begin{itemize}
			\item Ex.~7：从一维箱出发推到三维平动配分函数；
			\item Ex.~24：用 $q_{\mathrm{trans}}\propto m^{3/2}$ 处理熵比；
			\item Ex.~48：估算 0 到 $kT$ 之间有多少平动能级。
		\end{itemize}
	\end{examplebox}

	\section{双原子分子的转动配分函数}

	\subsection{刚性转子能级}
	双原子刚性转子：
	\[
		E_J=BJ(J+1),\qquad J=0,1,2,\dots
	\]
	简并度为
	\[
		g_J=2J+1.
	\]
	故
	\[
		q_{\mathrm{rot}}=\sum_{J=0}^{\infty}(2J+1)e^{-BJ(J+1)/kT}.
	\]
	前几项是
	\[
		q_{\mathrm{rot}}=1+3e^{-2B/kT}+5e^{-6B/kT}+\cdots.
	\]

	其中
	\[
		B=\frac{\hbar^2}{2I},
		\qquad
		I=\mu R^2,
		\qquad
		\mu=\frac{m_1m_2}{m_1+m_2}.
	\]
	若用波数表示：
	\[
		\tilde E_J=\tilde BJ(J+1),
		\qquad
		\tilde B=\frac{h}{8\pi^2 cI}.
	\]

	\begin{keypoint}
		转动常数最容易混两套写法：若写
		\[
			E_J=BJ(J+1),
		\]
		则 $B$ 是\textbf{能量}；若写
		\[
			\tilde E_J=\tilde B J(J+1),
		\]
		则 $\tilde B$ 是\textbf{波数}（常用 $\mathrm{cm^{-1}}$）。Boltzmann 因子里分子分母必须同单位：前者配能量形式的 $kT$，后者则可配波数形式的 $kT$，此时常用
		\[
			k=0.6950\ \mathrm{cm^{-1}\,K^{-1}}.
		\]
	\end{keypoint}

	\subsection{高温近似与转动温度}
	室温下常用换算
	\[
		k=0.6950\ \mathrm{cm^{-1}K^{-1}},
		\qquad
		kT\approx 207\ \mathrm{cm^{-1}}\ (298\,\mathrm K).
	\]
	对多数双原子分子，$kT\gg B$，因此可把和式近似为积分：
	\[
		q_{\mathrm{rot}}\approx \int_0^{\infty}(2J+1)e^{-BJ(J+1)/kT}\,dJ
		=\frac{kT}{B}.
	\]

	定义转动温度
	\[
		\theta_{\mathrm{rot}}=\frac{B}{k},
	\]
	再加上对称数 $\sigma$：
	\[
		q_{\mathrm{rot}}=\frac{kT}{\sigma B}=\frac{T}{\sigma\theta_{\mathrm{rot}}}.
	\]
	其中
	\[
		\sigma=1\ (\text{异核}),\qquad \sigma=2\ (\text{同核双原子}).
	\]

	\begin{keypoint}
		双原子高温转动公式：
		\[
			q_{\mathrm{rot}}=\frac{T}{\sigma\theta_{\mathrm{rot}}}.
		\]
		能不能直接用它，先看 $T\gg\theta_{\mathrm{rot}}$ 是否成立。
	\end{keypoint}

	\subsection{低温时必须逐项算}
	若 $T\not\gg\theta_{\mathrm{rot}}$，积分公式失效，只能回到
	\[
		q_{\mathrm{rot}}=\sum_{J=0}^{\infty}(2J+1)e^{-BJ(J+1)/kT}
	\]
	逐项累加到项足够小为止。此时只占据低几个转动态，所以其实不难算。

	\begin{examplebox}
		习题对应：
		\begin{itemize}
			\item Ex.~8：室温下 $N_2$ 的 $\theta_{\mathrm{rot}}$ 与 $q_{\mathrm{rot}}$；
			\item Ex.~21：30 K 下不能偷用高温公式，必须逐项求和；
			\item Ex.~47：用转动分布找最可几 $J$。
		\end{itemize}
	\end{examplebox}

	\section{非线性分子的转动配分函数}

	非线性分子有三个转动惯量 $I_a,I_b,I_c$，对应三个转动常数和三个转动温度。在高温极限，
	\[
		q_{\mathrm{rot}}=\frac{\pi^{1/2}}{\sigma}
		\left(\frac{T^3}{\theta_{\mathrm{rot},a}\theta_{\mathrm{rot},b}\theta_{\mathrm{rot},c}}\right)^{1/2}.
	\]
	典型例子：$H_2O$ 的 $\sigma=2$，$NH_3$ 的 $\sigma=3$。

	\begin{keypoint}
		线性分子：一个“转动温度”；
		非线性分子：三个“转动温度”。
		高温极限下仍是“很多转动态可及”的思想，只是公式变成三轴版本。
	\end{keypoint}

	\section{振动配分函数}

	\subsection{简谐振子模型}
	讲义把振动看作简谐振子，其能级（相对势阱底）为
	\[
		E_v=\left(v+\frac12\right)h\nu_0,
		\qquad v=0,1,2,\dots
	\]
	若把基态定义为零点，则
	\[
		E_v=vh\nu_0.
	\]
	于是
	\[
		q'_{\mathrm{vib}}=\sum_{v=0}^{\infty}e^{-vh\nu_0/kT}
		=1+e^{-h\nu_0/kT}+e^{-2h\nu_0/kT}+\cdots.
	\]

	这是几何级数，令
	\[
		\theta_{\mathrm{vib}}=\frac{h\nu_0}{k},
	\]
	则
	\[
		q'_{\mathrm{vib}}=\frac{1}{1-e^{-\theta_{\mathrm{vib}}/T}}.
	\]
	若能量相对势阱底而不是相对基态，则
	\[
		q_{\mathrm{vib}}=\frac{e^{-\theta_{\mathrm{vib}}/(2T)}}{1-e^{-\theta_{\mathrm{vib}}/T}}.
	\]

	\begin{keypoint}
		振动一定要分清零点：
		\[
			q_{\mathrm{vib}}=\frac{e^{-\theta_{\mathrm{vib}}/(2T)}}{1-e^{-\theta_{\mathrm{vib}}/T}}
			\quad (\text{相对势阱底}),
		\]
		\[
			q'_{\mathrm{vib}}=\frac{1}{1-e^{-\theta_{\mathrm{vib}}/T}}
			\quad (\text{相对基态}).
		\]
		其中 $\theta_{\mathrm{vib}}$ 的单位是 K；而 $h\nu_0$、$k\theta_{\mathrm{vib}}$、$hc\tilde\nu_0$ 是同一能级间隔的不同写法。做指数时只要保证分子分母同单位即可。
		习题 9 和 18 明确会考这两个版本。
	\end{keypoint}

	\subsection{高低温判断}
	振动能级间距就是 $h\nu_0$ 或 $k\theta_{\mathrm{vib}}$。通常分子振动频率较高，室温下常有
	\[
		\theta_{\mathrm{vib}}\gg T,
	\]
	故大多数振动模式接近“冻结”，只有基态显著占据，$q'_{\mathrm{vib}}$ 接近 1。只有在较高温或低频振动下，振动贡献才明显。

	\begin{examplebox}
		典型对应：
		\begin{itemize}
			\item Ex.~9：从几何级数求 $q_{\mathrm{vib}}$；
			\item Ex.~10：Morse 振子与 SHO 近似比较；
			\item Ex.~18/19：由 $q_{\mathrm{vib}}$ 求 $U,C_V$，分析高低温极限。
		\end{itemize}
	\end{examplebox}

	\section{本讲解题模板}

	\subsection{平动题}
	\begin{enumerate}
		\item 写粒子在箱中能级；
		\item 估计 $n_{\max}$，证明“很多项贡献”；
		\item 用积分得 $q_{\mathrm{trans},1D}$；
		\item 再乘到三维，写成 $V/\Lambda^3$。
	\end{enumerate}

	\subsection{转动题}
	\begin{enumerate}
		\item 先算 $\theta_{\mathrm{rot}}=B/k$；
		\item 判断 $T\gg\theta_{\mathrm{rot}}$ 还是 $T\sim\theta_{\mathrm{rot}}$；
		\item 高温就用 $T/(\sigma\theta_{\mathrm{rot}})$；
		\item 低温则回到和式逐项累加。
	\end{enumerate}

	\subsection{振动题}
	\begin{enumerate}
		\item 先搞清楚能量零点是在势阱底还是基态；
		\item 再写几何级数求和；
		\item 最后用 $\theta_{\mathrm{vib}}=h\nu_0/k$ 判断是否冻结。
	\end{enumerate}

	\section{最后速记}

	\begin{keypoint}
		Lect.~4 高频公式：
		\[
			q=\sum_j g_j e^{-\varepsilon_j/kT},
			\qquad
			q_{\mathrm{trans},3D}=\left(\frac{2\pi mkT}{h^2}\right)^{3/2}V=\frac{V}{\Lambda^3},
		\]
		\[
			\Lambda=\left(\frac{h^2}{2\pi mkT}\right)^{1/2},
			\qquad
			q_{\mathrm{rot}}=\sum_{J=0}^{\infty}(2J+1)e^{-BJ(J+1)/kT},
		\]
		\[
			q_{\mathrm{rot}}=\frac{T}{\sigma\theta_{\mathrm{rot}}}\ (T\gg\theta_{\mathrm{rot}}),
			\qquad
			\theta_{\mathrm{rot}}=\frac{B}{k},
		\]
		\[
			q_{\mathrm{rot,nonlinear}}=\frac{\pi^{1/2}}{\sigma}
			\left(\frac{T^3}{\theta_a\theta_b\theta_c}\right)^{1/2},
		\]
		\[
			q'_{\mathrm{vib}}=\frac{1}{1-e^{-\theta_{\mathrm{vib}}/T}},
			\qquad
			q_{\mathrm{vib}}=\frac{e^{-\theta_{\mathrm{vib}}/(2T)}}{1-e^{-\theta_{\mathrm{vib}}/T}}.
		\]
		\[
			T\not\gg \theta_{\mathrm{rot}}\Rightarrow q_{\mathrm{rot}}\text{ 必须回到求和式逐项算，}
			\qquad
			\theta_{\mathrm{vib}}\gg T\Rightarrow q'_{\mathrm{vib}}\approx 1.
		\]
	\end{keypoint}

\end{multicols}
\end{document}
